%% kapa: exam=1 round=31 year=2020 period=1 form=A subject=economics subject_ko=경제학원론 seq=28 number=68
%% source: Q-Net 감정평가사 기출문제 (cst00309) artlSeq=5207219 file=extracted/1008/5207219/A형/A형/2020년도 제31회 감정평가사 1차 시험_1교시_A형.pdf
%% conversion: kordoc 4.24.51 -> tools/gen1.py (auto); review=proofread-sonnet; flags=none
\begin{question}{68}
2015년과 2020년 빅맥 가격이 아래와 같다. 일물일가의 법칙이 성립할 때, 옳지 않은 것은? (단, 환율은 빅맥 가격을 기준으로 표시한다.) \par
{\small\centering
\begin{tabular}{|c|c|c|c|}\hline
\multicolumn{2}{|c|}{2015년} & \multicolumn{2}{c|}{2020년} \\ \hline
원화 가격 & 달러 가격 & 원화 가격 & 달러 가격 \\ \hline
5,000원 & 5달러 & 5,400원 & 6달러 \\ \hline
\end{tabular}\par}
\begin{choices}
  \item 빅맥의 원화 가격은 두 기간 사이에 8\% 상승했다.
  \item 빅맥의 1달러 당 원화 가격은 두 기간 사이에 10\% 하락했다.
  \item 달러 대비 원화의 가치는 두 기간 사이에 10\% 상승했다.
  \item 달러 대비 원화의 실질환율은 두 기간 사이에 변하지 않았다.
  \item 2020년 원화의 명목환율은 구매력평가 환율보다 낮다.
\end{choices}
\end{question}
\answer{5}
